\documentclass[a4paper,10pt]{article} \usepackage{anysize}
\marginsize{2cm}{2cm}{1cm}{1cm}
%\textwidth 6.0in \textheight = 664pt
\usepackage{xltxtra}
\usepackage{xunicode} \usepackage{graphicx}
\usepackage{color} \usepackage{xgreek} \usepackage{fancyvrb}
\usepackage{minted}
\usepackage{listings}
\usepackage{enumitem} \usepackage{framed} \usepackage{relsize}
\usepackage{float} \setmainfont[Mapping=TeX-text]{CMU Serif}
\begin{document}

\begin{titlepage}
\begin{center}
\begin{figure}[t] 
     \includegraphics[scale=0.7]{title/ntua_logo}
\end{figure}
\begin{LARGE}\textbf{ΕΘΝΙΚΟ ΜΕΤΣΟΒΙΟ ΠΟΛΥΤΕΧΝΕΙΟ\\}\end{LARGE}
\vspace{2cm}
\begin{Large}
ΣΧΟΛΗ ΗΜ\&ΜΥ\\
ΤΟΜΕΑΣ ΤΕΧΝΟΛΟΓΙΑΣ ΕΠΙΚΟΙΝΩΝΙΩΝ, ΗΛΕΚΤΡΟΝΙΚΗΣ ΚΑΙ ΣΥΣΤΗΜΑΤΩΝ ΠΛΗΡΟΦΟΡΙΚΗΣ\\
ΕΡΓΑΣΤΗΡΙΟ ΨΗΦΙΑΚΩΝ ΣΥΣΤΗΜΑΤΩΝ\\
1\textsuperscript{η} ΑΣΚΗΣΗ\\
Ακ. έτος 2010-2011\\
\end{Large}
\vspace{5cm}
\Large Τμήμα Γ, Ομάδα 2\textsuperscript{η}\\
\vspace{1cm}
\begin{tabular}{l r}
\Large{Γερακάρης Βασίλης}&
\large{Α.Μ.: 03108137}\\
\Large{Λύρας Γρηγόρης}&
\large{Α.Μ.: 03109687}\\
\end{tabular}\\
\vspace{5cm}

\vfill
\large\today\\
\end{center}
\end{titlepage}



\section*{} \setcounter{section}{1}
\subsection{Δημιουργία δεδομένου δέντρου διεργασιών} Ο πηγαίος κώδικας της main.c που
κληθήκαμε να γράψουμε ήταν ο εξής:

\inputminted[linenos,fontsize=\footnotesize,frame=leftline]{c}{files/ask2-fork.c}
\emph{Ερωτήσεις}
\begin{enumerate}
\item Αν η διεργασία A τερματίσει πρόωρα τα παιδιά της (B,C) θα συνεχίσουν να
εκτελούνται και υιοθετούνται από την INIT. Όταν αυτά τερματίσουν η INIT θα
κάνει wait και θα φύγουν από τη μνήμη.
\item Φαίνεται μία ακόμη διεργασία. Αυτή είναι η διεργασία που ξεκίνησε τα
forks δημιουργώντας την A.
\item Σε διαφορετική περίπτωση ένας χρήστης θα μπορούσε να εκτελέσει
απεριόριστα διεργασίες. Κάτι τέτοιο δεν επιτρέπεται μιας και η ανεξέλεγκτη
κατανάλωση πόρων συστήματος μπορεί να καταστήσει το σύστημά μας μη
λειτουργικό, κάτι που ο διαχειριστής πρέπει να αποφύγει.
\end{enumerate}


\subsection{Δημιουργία αυθαίρετου δέντρου διεργασιών} Ο πηγαίος κώδικας της main.c που
\inputminted[linenos,fontsize=\footnotesize,frame=leftline]{c}{files/ask2-tree.c}
\emph{Ερωτήσεις}
\begin{enumerate}
\item Οι διεργασίες γεννιούνται κατά επίπεδο. Παρόλα αυτά, δεν είναι
εξασφαλισμένη η σειρά που θα εκτελεστούν οι διεργασίες. Έχει να κάνει με τη
χρονοδρομολόγηση (δηλαδή το D του προηγούμενου ερωτήματος θα μπορούσε να
εκτελεστεί πριν από το C αν η χρονοδρομολόγηση το ευνοούσε. 
Ωστόσο το κβάντο χρόνου είναι τέτοιο η εμφάνιση των
μηνυμάτων δημιουργίας μοιάζει σχεδόν σειριακή.
\end{enumerate}


\subsection{Αποστολή και χειρισμός σημάτων}

\inputminted[linenos,fontsize=\footnotesize,frame=leftline]{c}{files/ask2-signals.c}
\emph{Ερωτήσεις}
\begin{enumerate}
\item Είναι σαφώς καύτερη και πιο αξιόπιστη καθώς δεν στηρίζεται στη
χρονοδρομολόγηση των διεργασιών.
\item Προτού στείλουμε σήματα στα παιδιά πρέπει να έχουμε βεβαιωθεί πως αυτά
είναι σε ready state ώστε να δεχτούν σήμα. Αν το κάνουμε νωρίτερα αυτό θα
αγνοηθεί. Συνεπώς περιμένουμε μέχρι όλα τα παιδιά να κάνουν raise(SIGSTOP)
ώστε να δεχτούν το SIGCONT.
\end{enumerate}

\subsection{Παράλληλος υπολογισμός αριθμητικής έφρασης}
\inputminted[linenos,fontsize=\footnotesize,frame=leftline]{c}{files/ask2-pipes.c}
\end{document}
